\IfFileExists{stacks-project.cls}{%
\documentclass{stacks-project}
}{%
\documentclass{amsart}
}

% For dealing with references we use the comment environment
\usepackage{verbatim}
\newenvironment{reference}{\comment}{\endcomment}
%\newenvironment{reference}{}{}
\newenvironment{slogan}{\comment}{\endcomment}
\newenvironment{history}{\comment}{\endcomment}

% For commutative diagrams we use Xy-pic
\usepackage[all]{xy}

% We use 2cell for 2-commutative diagrams.
\xyoption{2cell}
\UseAllTwocells

% We use multicol for the list of chapters between chapters
\usepackage{multicol}

% This is generally recommended for better output
\usepackage{lmodern}
\usepackage[T1]{fontenc}

% For cross-file-references
\usepackage{xr-hyper}

% Package for hypertext links:
\usepackage{hyperref}

% For any local file, say "hello.tex" you want to link to please
% use \externaldocument[hello-]{hello}
\externaldocument[introduction-]{introduction}
\externaldocument[conventions-]{conventions}
\externaldocument[sets-]{sets}
\externaldocument[categories-]{categories}
\externaldocument[topology-]{topology}
\externaldocument[sheaves-]{sheaves}
\externaldocument[sites-]{sites}
\externaldocument[stacks-]{stacks}
\externaldocument[fields-]{fields}
\externaldocument[algebra-]{algebra}
\externaldocument[brauer-]{brauer}
\externaldocument[homology-]{homology}
\externaldocument[derived-]{derived}
\externaldocument[simplicial-]{simplicial}
\externaldocument[more-algebra-]{more-algebra}
\externaldocument[smoothing-]{smoothing}
\externaldocument[modules-]{modules}
\externaldocument[sites-modules-]{sites-modules}
\externaldocument[injectives-]{injectives}
\externaldocument[cohomology-]{cohomology}
\externaldocument[sites-cohomology-]{sites-cohomology}
\externaldocument[dga-]{dga}
\externaldocument[dpa-]{dpa}
\externaldocument[sdga-]{sdga}
\externaldocument[hypercovering-]{hypercovering}
\externaldocument[schemes-]{schemes}
\externaldocument[constructions-]{constructions}
\externaldocument[properties-]{properties}
\externaldocument[morphisms-]{morphisms}
\externaldocument[coherent-]{coherent}
\externaldocument[divisors-]{divisors}
\externaldocument[limits-]{limits}
\externaldocument[varieties-]{varieties}
\externaldocument[topologies-]{topologies}
\externaldocument[descent-]{descent}
\externaldocument[perfect-]{perfect}
\externaldocument[more-morphisms-]{more-morphisms}
\externaldocument[flat-]{flat}
\externaldocument[groupoids-]{groupoids}
\externaldocument[more-groupoids-]{more-groupoids}
\externaldocument[etale-]{etale}
\externaldocument[chow-]{chow}
\externaldocument[intersection-]{intersection}
\externaldocument[pic-]{pic}
\externaldocument[weil-]{weil}
\externaldocument[adequate-]{adequate}
\externaldocument[dualizing-]{dualizing}
\externaldocument[duality-]{duality}
\externaldocument[discriminant-]{discriminant}
\externaldocument[derham-]{derham}
\externaldocument[local-cohomology-]{local-cohomology}
\externaldocument[algebraization-]{algebraization}
\externaldocument[curves-]{curves}
\externaldocument[resolve-]{resolve}
\externaldocument[models-]{models}
\externaldocument[functors-]{functors}
\externaldocument[equiv-]{equiv}
\externaldocument[pione-]{pione}
\externaldocument[etale-cohomology-]{etale-cohomology}
\externaldocument[proetale-]{proetale}
\externaldocument[relative-cycles-]{relative-cycles}
\externaldocument[more-etale-]{more-etale}
\externaldocument[trace-]{trace}
\externaldocument[crystalline-]{crystalline}
\externaldocument[spaces-]{spaces}
\externaldocument[spaces-properties-]{spaces-properties}
\externaldocument[spaces-morphisms-]{spaces-morphisms}
\externaldocument[decent-spaces-]{decent-spaces}
\externaldocument[spaces-cohomology-]{spaces-cohomology}
\externaldocument[spaces-limits-]{spaces-limits}
\externaldocument[spaces-divisors-]{spaces-divisors}
\externaldocument[spaces-over-fields-]{spaces-over-fields}
\externaldocument[spaces-topologies-]{spaces-topologies}
\externaldocument[spaces-descent-]{spaces-descent}
\externaldocument[spaces-perfect-]{spaces-perfect}
\externaldocument[spaces-more-morphisms-]{spaces-more-morphisms}
\externaldocument[spaces-flat-]{spaces-flat}
\externaldocument[spaces-groupoids-]{spaces-groupoids}
\externaldocument[spaces-more-groupoids-]{spaces-more-groupoids}
\externaldocument[bootstrap-]{bootstrap}
\externaldocument[spaces-pushouts-]{spaces-pushouts}
\externaldocument[spaces-chow-]{spaces-chow}
\externaldocument[groupoids-quotients-]{groupoids-quotients}
\externaldocument[spaces-more-cohomology-]{spaces-more-cohomology}
\externaldocument[spaces-simplicial-]{spaces-simplicial}
\externaldocument[spaces-duality-]{spaces-duality}
\externaldocument[formal-spaces-]{formal-spaces}
\externaldocument[restricted-]{restricted}
\externaldocument[spaces-resolve-]{spaces-resolve}
\externaldocument[formal-defos-]{formal-defos}
\externaldocument[defos-]{defos}
\externaldocument[cotangent-]{cotangent}
\externaldocument[examples-defos-]{examples-defos}
\externaldocument[algebraic-]{algebraic}
\externaldocument[examples-stacks-]{examples-stacks}
\externaldocument[stacks-sheaves-]{stacks-sheaves}
\externaldocument[criteria-]{criteria}
\externaldocument[artin-]{artin}
\externaldocument[quot-]{quot}
\externaldocument[stacks-properties-]{stacks-properties}
\externaldocument[stacks-morphisms-]{stacks-morphisms}
\externaldocument[stacks-limits-]{stacks-limits}
\externaldocument[stacks-cohomology-]{stacks-cohomology}
\externaldocument[stacks-perfect-]{stacks-perfect}
\externaldocument[stacks-introduction-]{stacks-introduction}
\externaldocument[stacks-more-morphisms-]{stacks-more-morphisms}
\externaldocument[stacks-geometry-]{stacks-geometry}
\externaldocument[moduli-]{moduli}
\externaldocument[moduli-curves-]{moduli-curves}
\externaldocument[examples-]{examples}
\externaldocument[exercises-]{exercises}
\externaldocument[guide-]{guide}
\externaldocument[desirables-]{desirables}
\externaldocument[coding-]{coding}
\externaldocument[obsolete-]{obsolete}
\externaldocument[fdl-]{fdl}
\externaldocument[index-]{index}

% Theorem environments.
%
\theoremstyle{plain}
\newtheorem{theorem}[subsection]{Theorem}
\newtheorem{proposition}[subsection]{Proposition}
\newtheorem{lemma}[subsection]{Lemma}

\theoremstyle{definition}
\newtheorem{definition}[subsection]{Definition}
\newtheorem{example}[subsection]{Example}
\newtheorem{exercise}[subsection]{Exercise}
\newtheorem{situation}[subsection]{Situation}

\theoremstyle{remark}
\newtheorem{remark}[subsection]{Remark}
\newtheorem{remarks}[subsection]{Remarks}

\numberwithin{equation}{subsection}

% Macros
%
\def\lim{\mathop{\mathrm{lim}}\nolimits}
\def\colim{\mathop{\mathrm{colim}}\nolimits}
\def\Spec{\mathop{\mathrm{Spec}}}
\def\Hom{\mathop{\mathrm{Hom}}\nolimits}
\def\Ext{\mathop{\mathrm{Ext}}\nolimits}
\def\SheafHom{\mathop{\mathcal{H}\!\mathit{om}}\nolimits}
\def\SheafExt{\mathop{\mathcal{E}\!\mathit{xt}}\nolimits}
\def\Sch{\mathit{Sch}}
\def\Mor{\mathop{\mathrm{Mor}}\nolimits}
\def\Ob{\mathop{\mathrm{Ob}}\nolimits}
\def\Sh{\mathop{\mathit{Sh}}\nolimits}
\def\NL{\mathop{N\!L}\nolimits}
\def\CH{\mathop{\mathrm{CH}}\nolimits}
\def\proetale{{pro\text{-}\acute{e}tale}}
\def\etale{{\acute{e}tale}}
\def\QCoh{\mathit{QCoh}}
\def\Ker{\mathop{\mathrm{Ker}}}
\def\Im{\mathop{\mathrm{Im}}}
\def\Coker{\mathop{\mathrm{Coker}}}
\def\Coim{\mathop{\mathrm{Coim}}}

% Boxtimes
%
\DeclareMathSymbol{\boxtimes}{\mathbin}{AMSa}{"02}

%
% Macros for moduli stacks/spaces
%
\def\QCohstack{\mathcal{QC}\!\mathit{oh}}
\def\Cohstack{\mathcal{C}\!\mathit{oh}}
\def\Spacesstack{\mathcal{S}\!\mathit{paces}}
\def\Quotfunctor{\mathrm{Quot}}
\def\Hilbfunctor{\mathrm{Hilb}}
\def\Curvesstack{\mathcal{C}\!\mathit{urves}}
\def\Polarizedstack{\mathcal{P}\!\mathit{olarized}}
\def\Complexesstack{\mathcal{C}\!\mathit{omplexes}}
% \Pic is the operator that assigns to X its picard group, usage \Pic(X)
% \Picardstack_{X/B} denotes the Picard stack of X over B
% \Picardfunctor_{X/B} denotes the Picard functor of X over B
\def\Pic{\mathop{\mathrm{Pic}}\nolimits}
\def\Picardstack{\mathcal{P}\!\mathit{ic}}
\def\Picardfunctor{\mathrm{Pic}}
\def\Deformationcategory{\mathcal{D}\!\mathit{ef}}
\usepackage{cleveref}
\crefname{originaldefinition}{original definition}{original definitions}
\Crefname{originaldefinition}{Original definition}{Original definitions}
\makeatletter
\newlabel{stacks-definition-topology-inherited}{{06NV}{source}{}{sourcepointer-1}{}}
\newlabel{stacks-definition-topology-inherited@cref}{{[originaldefinition][0][]06NV}{[1][0][]source}{}{}{}}
\makeatother
\crefname{originaldefinition}{original definition}{original definitions}
\Crefname{originaldefinition}{Original definition}{Original definitions}
\makeatletter
\newlabel{stacks-definition-stack-in-groupoids}{{02ZI}{source}{}{sourcepointer-2}{}}
\newlabel{stacks-definition-stack-in-groupoids@cref}{{[originaldefinition][0][]02ZI}{[1][0][]source}{}{}{}}
\makeatother
\crefname{originallemma}{original lemma}{original lemmas}
\Crefname{originallemma}{Original lemma}{Original lemmas}
\makeatletter
\newlabel{lemma-polarized-fibred-in-groupoids}{{0D3Z}{source}{}{sourcepointer-3}{}}
\newlabel{lemma-polarized-fibred-in-groupoids@cref}{{[originallemma][0][]0D3Z}{[1][0][]source}{}{}{}}
\makeatother
\crefname{originalproposition}{original proposition}{original propositions}
\Crefname{originalproposition}{Original proposition}{Original propositions}
\makeatletter
\newlabel{spaces-descent-proposition-fpqc-descent-quasi-coherent}{{04W8}{source}{}{sourcepointer-4}{}}
\newlabel{spaces-descent-proposition-fpqc-descent-quasi-coherent@cref}{{[originalproposition][0][]04W8}{[1][0][]source}{}{}{}}
\makeatother
\crefname{originallemma}{original lemma}{original lemmas}
\Crefname{originallemma}{Original lemma}{Original lemmas}
\makeatletter
\newlabel{spaces-descent-lemma-descending-property-ample}{{0D3C}{source}{}{sourcepointer-5}{}}
\newlabel{spaces-descent-lemma-descending-property-ample@cref}{{[originallemma][0][]0D3C}{[1][0][]source}{}{}{}}
\makeatother
\crefname{originallemma}{original lemma}{original lemmas}
\Crefname{originallemma}{Original lemma}{Original lemmas}
\makeatletter
\newlabel{lemma-spaces-stack}{{0D1G}{source}{}{sourcepointer-6}{}}
\newlabel{lemma-spaces-stack@cref}{{[originallemma][0][]0D1G}{[1][0][]source}{}{}{}}
\makeatother
\crefname{originallemma}{original lemma}{original lemmas}
\Crefname{originallemma}{Original lemma}{Original lemmas}
\makeatletter
\newlabel{stacks-lemma-stack-over-stack}{{09WX}{source}{}{sourcepointer-7}{}}
\newlabel{stacks-lemma-stack-over-stack@cref}{{[originallemma][0][]09WX}{[1][0][]source}{}{}{}}
\makeatother
\title{The polarized stack is a stack in groupoids}
\author{The Stacks Project Authors}
\date{stacks-project:quot.tex:local-snapshot-20261012; GNU FDL 1.2 or later}
\begin{document}
\maketitle
\setcounter{subsection}{1}
\noindent\textit{Editorial note.} All delivered mathematical text is the author's own, copied from the pinned Stacks Project chapter file quot.tex; no statement or proof step was rewritten, reordered or removed. The counted span is the author environment and its own proof environment. References inside the span are handed to the assembler as navigation stubs that carry the canonical tag of the cited tag only; whether each of them is load-bearing is recorded on the stub itself.\par\medskip
\noindent\textit{Allowed standard prerequisites (editorial).}\begin{itemize}
\item Definition stacks-definition-topology-inherited (canonical tag 06NV). cited at statement line 4087. Inherited topology on the base category.
\item Definition stacks-definition-stack-in-groupoids (canonical tag 02ZI). cited at proof line 4095. The three stack conditions.
\item Lemma lemma-polarized-fibred-in-groupoids (canonical tag 0D3Z). cited at proof lines 4096-4097. Condition (1) already known.
\item Proposition spaces-descent-proposition-fpqc-descent-quasi-coherent (canonical tag 04W8). cited at proof lines 4127-4128 and 4151-4152. Descent of quasi-coherent modules gives unique gluings.
\item Lemma spaces-descent-lemma-descending-property-ample (canonical tag 0D3C). cited at proof lines 4160-4161. Ampleness descends.
\item Lemma lemma-spaces-stack (canonical tag 0D1G). cited at proof lines 4164-4165. The base is a stack in groupoids over Sch\_fppf.
\item Lemma stacks-lemma-stack-over-stack (canonical tag 09WX). cited at proof lines 4166-4167. Formal transfer to Sch\_fppf.
\end{itemize}\medskip
\ifdefined\hypertarget\hypertarget{sourcepointer-1}{}\fi
\noindent\textit{Source reference.} Original definition 06NV: sources/stacks-chow/tags.txt, source lines 8549--8549. Cited at statement line 4087: the topology on Spaces'\_\{fp, flat, proper\} is the inherited one; load-bearing for the first assertion.\par\medskip
\ifdefined\hypertarget\hypertarget{sourcepointer-2}{}\fi
\noindent\textit{Source reference.} Original definition 02ZI: sources/stacks-chow/tags.txt, source lines 3913--3913. Cited at proof line 4095: conditions (1), (2), (3) being checked; load-bearing.\par\medskip
\ifdefined\hypertarget\hypertarget{sourcepointer-3}{}\fi
\noindent\textit{Source reference.} Original lemma 0D3Z: sources/stacks-chow/tags.txt, source lines 16434--16434. Cited at proof lines 4096-4097: condition (1) already verified; load-bearing.\par\medskip
\ifdefined\hypertarget\hypertarget{sourcepointer-4}{}\fi
\noindent\textit{Source reference.} Original proposition 04W8: sources/stacks-chow/tags.txt, source lines 6351--6351. Cited at proof lines 4127-4128 (descent for morphisms) and 4151-4152 (descent for objects): unique gluing of isomorphisms and of the invertible module; load-bearing twice.\par\medskip
\ifdefined\hypertarget\hypertarget{sourcepointer-5}{}\fi
\noindent\textit{Source reference.} Original lemma 0D3C: sources/stacks-chow/tags.txt, source lines 16412--16412. Cited at proof lines 4160-4161: the descended L is ample; load-bearing.\par\medskip
\ifdefined\hypertarget\hypertarget{sourcepointer-6}{}\fi
\noindent\textit{Source reference.} Original lemma 0D1G: sources/stacks-chow/tags.txt, source lines 16346--16346. Cited at proof lines 4164-4165: base stack is a stack in groupoids; load-bearing.\par\medskip
\ifdefined\hypertarget\hypertarget{sourcepointer-7}{}\fi
\noindent\textit{Source reference.} Original lemma 09WX: sources/stacks-chow/tags.txt, source lines 12510--12510. Cited at proof lines 4166-4167: formal transfer of the stack property to Sch\_fppf; load-bearing.\par\medskip
\begin{lemma}
\label{lemma-polarized-stack}
The category $\Polarizedstack$ is a stack in groupoids over
$\Spacesstack'_{fp, flat, proper}$ (endowed with the inherited topology,
see Stacks, Definition \ref{stacks-definition-topology-inherited}).
The category $\Polarizedstack$ is a stack in groupoids over $\Sch_{fppf}$.
\end{lemma}

\begin{proof}
We prove $\Polarizedstack$ is a stack in groupoids over
$\Spacesstack'_{fp, flat, proper}$
by checking conditions (1), (2), and (3)
of Stacks, Definition \ref{stacks-definition-stack-in-groupoids}.
We have already seen (1) in
Lemma \ref{lemma-polarized-fibred-in-groupoids}.

\medskip\noindent
A covering of $\Spacesstack'_{fp, flat, proper}$ comes about
in the following manner: Let $X \to S$ be an object of
$\Spacesstack'_{fp, flat, proper}$. Suppose that
$\{S_i \to S\}_{i \in I}$ is a covering of $\Sch_{fppf}$.
Set $X_i = S_i \times_S X$. Then $\{(X_i \to S_i) \to (X \to S)\}_{i \in I}$
is a covering of $\Spacesstack'_{fp, flat, proper}$ and
every covering of $\Spacesstack'_{fp, flat, proper}$ is isomorphic
to one of these. Set $S_{ij} = S_i \times_S S_j$ and
$X_{ij} = S_{ij} \times_S X$ so that $(X_{ij} \to S_{ij}) =
(X_i \to S_i) \times_{(X \to S)} (X_j \to S_j)$.
Next, suppose that $\mathcal{L}, \mathcal{N}$
are ample invertible sheaves on $X/S$ so that
$(X \to S, \mathcal{L})$ and $(X \to S, \mathcal{N})$
are two objects of $\Polarizedstack$ over the object $(X \to S)$.
To check descent for morphisms, we assume we have morphisms
$(\text{id}, \text{id}, \varphi_i)$ from
$(X_i \to S_i, \mathcal{L}|_{X_i})$ to
$(X_i \to S_i, \mathcal{N}|_{X_i})$
whose base changes to morphisms from
$(X_{ij} \to S_{ij}, \mathcal{L}|_{X_{ij}})$ to
$(X_{ij} \to S_{ij}, \mathcal{N}|_{X_{ij}})$
agree. Then
$\varphi_i : \mathcal{L}|_{X_i} \to \mathcal{N}|_{X_i}$
are isomorphisms of invertible modules over $X_i$ such that
$\varphi_i$ and $\varphi_j$ restrict to the same
isomorphisms over $X_{ij}$.
By descent for quasi-coherent sheaves
(Descent on Spaces, Proposition
\ref{spaces-descent-proposition-fpqc-descent-quasi-coherent})
we obtain a unique isomorphism $\varphi : \mathcal{L} \to \mathcal{N}$
whose restriction to $X_i$ recovers $\varphi_i$.

\medskip\noindent
Decent for objects is proved in exactly the same manner.
Namely, suppose that
$\{(X_i \to S_i) \to (X \to S)\}_{i \in I}$
is a covering of $\Spacesstack'_{fp, flat, proper}$
as above.
Suppose we have objects $(X_i \to S_i, \mathcal{L}_i)$
of $\Polarizedstack$ lying over $(X_i \to S_i)$
and a descent datum
$$
(\text{id}, \text{id}, \varphi_{ij}) :
(X_{ij} \to S_{ij}, \mathcal{L}_i|_{X_{ij}})
\to
(X_{ij} \to S_{ij}, \mathcal{L}_j|_{X_{ij}})
$$
satisfying the obvious cocycle condition over
$(X_{ijk} \to S_{ijk})$ for every triple of indices.
Then by
 descent for quasi-coherent sheaves
(Descent on Spaces, Proposition
\ref{spaces-descent-proposition-fpqc-descent-quasi-coherent})
we obtain a unique invertible $\mathcal{O}_X$-module
$\mathcal{L}$ and isomorphisms $\mathcal{L}|_{X_i} \to \mathcal{L}_i$
recovering the descent datum $\varphi_{ij}$.
To show that
$(X \to S, \mathcal{L})$ is an object of
$\Polarizedstack$ we have to prove that
$\mathcal{L}$ is ample. This follows from
Descent on Spaces, Lemma \ref{spaces-descent-lemma-descending-property-ample}.

\medskip\noindent
Since we already have seen that $\Spacesstack'_{fp, flat, proper}$
is a stack in groupoids over $\Sch_{fppf}$
(Lemma \ref{lemma-spaces-stack}) it now follows formally
that $\Polarizedstack$ is a stack in groupoids over $\Sch_{fppf}$.
See Stacks, Lemma \ref{stacks-lemma-stack-over-stack}.
\end{proof}

\end{document}
